% ==============================================================================
% preamble/headings.tex — การจัดรูปแบบหัวข้อและการควบคุมหน้าขึ้นบทใหม่
% ตามมาตรฐาน thai-academic-latex-formatting
% ==============================================================================

% 1. รูปแบบบท (Chapter Formatting)
% - เริ่มต้นที่หน้าคี่ (Odd page / Right-hand page)
% - จัดชิดขวา (Right-aligned)
% - ป้ายชื่อบทเด่นชัด (Visually prominent chapter label)
\titleformat{\chapter}[display]
  {\normalfont\raggedleft}
  {\fontsize{18}{22}\selectfont\bfseries บทที่ \thechapter}
  {8pt}
  {\fontsize{20}{24}\selectfont\bfseries}

\titlespacing*{\chapter}{0pt}{30pt}{25pt}

% 2. รูปแบบหัวข้อย่อยระดับ 1 (Section)
% ป้องกันหัวข้อตกค้างท้ายหน้าด้วย \needspace
\newcommand{\academicsection}[1]{%
  \needspace{5\baselineskip}%
  \section{#1}%
}

\titleformat{\section}
  {\normalfont\fontsize{13}{16}\selectfont\bfseries}
  {\thesection}
  {1em}
  {}
\titlespacing*{\section}{0pt}{14pt plus 2pt minus 2pt}{6pt plus 1pt minus 1pt}

% 3. รูปแบบหัวข้อย่อยระดับ 2 (Subsection)
\newcommand{\academicsubsection}[1]{%
  \needspace{4\baselineskip}%
  \subsection{#1}%
}

\titleformat{\subsection}
  {\normalfont\fontsize{11.5}{14}\selectfont\bfseries}
  {\thesubsection}
  {0.8em}
  {}
\titlespacing*{\subsection}{0pt}{10pt plus 2pt minus 2pt}{4pt plus 1pt minus 1pt}

% 4. การตั้งค่าหัวกระดาษและท้ายกระดาษ (Running Headers & Footers)
\pagestyle{fancy}
\fancyhf{}
\renewcommand{\headrulewidth}{0.4pt}

% หน้าคู่ (Even): เลขหน้าชิดซ้าย, ชื่อหนังสือชิดขวา
\fancyhead[LE]{\fontsize{9}{11}\selectfont\thepage}
\fancyhead[RE]{\fontsize{9}{11}\selectfont คู่มือปฏิบัติการ Digital Agri xAI}

% หน้าคี่ (Odd): ชื่อบทชิดซ้าย, เลขหน้าชิดขวา
\fancyhead[LO]{\fontsize{9}{11}\selectfont\leftmark}
\fancyhead[RO]{\fontsize{9}{11}\selectfont\thepage}

% หน้าเริ่มต้นบท (Plain style) ให้มีเฉพาะเลขหน้าด้านล่างตรงกลาง
\fancypagestyle{plain}{
    \fancyhf{}
    \fancyfoot[C]{\fontsize{9}{11}\selectfont\thepage}
    \renewcommand{\headrulewidth}{0pt}
}

% ปรับแต่ง \chaptermark ให้แสดงชื่อบทในหัวกระดาษ
\renewcommand{\chaptermark}[1]{\markboth{บทที่ \thechapter\ #1}{}}
